% GENERATED FILE. Edit canonical lessons, then run npm run generate:books.
% canonical-source-hash: fnv1a64:be6f613a4d8d81c4
% canonical-lessons: FR-C217-presque, FR-C217-environ, FR-C217-seulement, FR-C217-ensemble, FR-C217-cependant

\chapter{Almost, About, Only, Together, Nevertheless}
\label{ch:fr-almost-about-only-together-nevertheless}
\hlchaptermodality{217}

\begin{chapteropening}
\textbf{By the end of this chapter:} \emph{I can say almost, about, only, together and nevertheless.}

Complete the last lesson of chapter 217: I can say almost, about, only, together and nevertheless.
\end{chapteropening}

\section[presque]{presque — almost}

\label{lesson:FR-C217-presque}



\begin{quote}
Before the new one: say the French for up there, then the French for the suburbs.
\end{quote}

\subsection*{You'll want to know: presque}
\textbf{presque} — \textquotedblleft{}almost\textquotedblright{}.

\begin{cousinweb}
\textbf{presque} — \emph{près}, \textquotedblleft{}near\textquotedblright{}, and \emph{que}, \textquotedblleft{}that\textquotedblright{}: nearly. \emph{Près} $\leftarrow$ Latin \emph{presse}, \textquotedblleft{}closely\textquotedblright{} (EtymDB 2.1).
\end{cousinweb}

\begin{grammarlens}[title={What you've built}]
One more word for this chapter.
\end{grammarlens}

\subsection*{Your turn}
\begin{itemize}
\raggedright
  \item \emph{Say it:} \emph{presque}
  \item \emph{Say it:} \emph{presque}, once more
  \item \emph{Say it:} say \emph{presque}
  \item \emph{Recall:} say the French for up there, then the French for the suburbs, then say \emph{presque} again
\end{itemize}

\begin{tcolorbox}[breakable,colback=teal!4,colframe=teal!35!black,title={Before you move on}]
What does presque mean? (\textquotedblleft{}Almost\textquotedblright{}.) Say it once more.
\end{tcolorbox}
\section[environ]{environ — about, approximately}

\label{lesson:FR-C217-environ}



\begin{quote}
Before the new one: say the French for the suburbs, then the French for almost.
\end{quote}

\subsection*{You'll want to know: environ}
\textbf{environ} — \textquotedblleft{}about, approximately\textquotedblright{}.

\begin{cousinweb}
\textbf{environ} — \emph{en}, \textquotedblleft{}in\textquotedblright{}, and Old French \emph{viron}, \textquotedblleft{}a circuit\textquotedblright{}, from the verb for \textquotedblleft{}to turn\textquotedblright{}: in the round of. \emph{En} $\leftarrow$ Latin \emph{in} (EtymDB 2.1); the verb behind \emph{viron} is of uncertain origin.
\end{cousinweb}

\begin{grammarlens}[title={What you've built}]
One more word for this chapter.
\end{grammarlens}

\subsection*{Your turn}
\begin{itemize}
\raggedright
  \item \emph{Say it:} \emph{environ}
  \item \emph{Say it:} \emph{environ}, once more
  \item \emph{Say it:} say \emph{environ}
  \item \emph{Recall:} say the French for the suburbs, then the French for almost, then say \emph{environ} again
\end{itemize}

\begin{tcolorbox}[breakable,colback=teal!4,colframe=teal!35!black,title={Before you move on}]
What does environ mean? (\textquotedblleft{}About, approximately\textquotedblright{}.) Say it once more.
\end{tcolorbox}
\section[seulement]{seulement — only}

\label{lesson:FR-C217-seulement}



\begin{quote}
Before the new one: say the French for almost, then the French for about.
\end{quote}

\subsection*{You'll want to know: seulement}
\textbf{seulement} — \textquotedblleft{}only\textquotedblright{}.

\begin{cousinweb}
\textbf{seulement} — the word for \textquotedblleft{}alone\textquotedblright{}, $\leftarrow$ Latin \emph{solus} (Lewis \& Short), with \emph{-ment}.
\end{cousinweb}

\begin{grammarlens}[title={What you've built}]
One more word for this chapter.
\end{grammarlens}

\subsection*{Your turn}
\begin{itemize}
\raggedright
  \item \emph{Say it:} \emph{seulement}
  \item \emph{Say it:} \emph{seulement}, once more
  \item \emph{Say it:} say \emph{seulement}
  \item \emph{Recall:} say the French for almost, then the French for about, then say \emph{seulement} again
\end{itemize}

\begin{tcolorbox}[breakable,colback=teal!4,colframe=teal!35!black,title={Before you move on}]
What does seulement mean? (\textquotedblleft{}Only\textquotedblright{}.) Say it once more.
\end{tcolorbox}
\section[ensemble]{ensemble — together}

\label{lesson:FR-C217-ensemble}



\begin{quote}
Before the new one: say the French for about, then the French for only.
\end{quote}

\subsection*{You'll want to know: ensemble}
\textbf{ensemble} — \textquotedblleft{}together\textquotedblright{}.

\begin{cousinweb}
\textbf{ensemble} $\leftarrow$ Late Latin \emph{insimul}, \textquotedblleft{}at the same time\textquotedblright{}, from \emph{simul}, \textquotedblleft{}together\textquotedblright{} — inherited. (EtymDB 2.1)
\end{cousinweb}

\begin{grammarlens}[title={What you've built}]
One more word for this chapter.
\end{grammarlens}

\subsection*{Your turn}
\begin{itemize}
\raggedright
  \item \emph{Say it:} \emph{ensemble}
  \item \emph{Say it:} \emph{ensemble}, once more
  \item \emph{Say it:} say \emph{ensemble}
  \item \emph{Recall:} say the French for about, then the French for only, then say \emph{ensemble} again
\end{itemize}

\begin{tcolorbox}[breakable,colback=teal!4,colframe=teal!35!black,title={Before you move on}]
What does ensemble mean? (\textquotedblleft{}Together\textquotedblright{}.) Say it once more.
\end{tcolorbox}
\section[cependant]{cependant — nevertheless, however}

\label{lesson:FR-C217-cependant}



\begin{quote}
Before the new one: say the French for only, then the French for together.
\end{quote}

\subsection*{You'll want to know: cependant}
\textbf{cependant} — \textquotedblleft{}nevertheless, however\textquotedblright{}.

\begin{cousinweb}
\textbf{cependant} — \emph{ce}, \textquotedblleft{}this\textquotedblright{}, and the word for \textquotedblleft{}during\textquotedblright{}, from the verb for \textquotedblleft{}to hang\textquotedblright{}, $\leftarrow$ Latin \emph{pendere} (EtymDB 2.1): \textquotedblleft{}during this\textquotedblright{}.
\end{cousinweb}

\begin{grammarlens}[title={What you've built}]
One more word for this chapter.
\end{grammarlens}

\subsection*{Your turn}
\begin{itemize}
\raggedright
  \item \emph{Say it:} \emph{cependant}
  \item \emph{Say it:} \emph{cependant}, once more
  \item \emph{Say it:} say \emph{cependant}
  \item \emph{Recall:} say the French for only, then the French for together, then say \emph{cependant} again
\end{itemize}

\begin{tcolorbox}[breakable,colback=teal!4,colframe=teal!35!black,title={Before you move on}]
What does cependant mean? (\textquotedblleft{}Nevertheless, however\textquotedblright{}.) Say it once more.
\end{tcolorbox}
